\documentclass[10pt]{article}

\linespread{1.1}
\usepackage[a4paper,margin=0.8in]{geometry}
\usepackage[T1]{fontenc}
\usepackage{lmodern}
\usepackage{amsmath,amssymb,mathtools}
\usepackage{array}
\usepackage{booktabs}
\usepackage{hyperref}
\usepackage{listings}

\hypersetup{
  colorlinks=true,
  linkcolor=blue,
  urlcolor=blue,
  citecolor=blue
}

\lstset{
  basicstyle=\ttfamily\small,
  columns=fullflexible,
  keepspaces=true,
  frame=single,
  showstringspaces=false
}

\newcommand{\Tr}{\mathrm{tr}}
\newcommand{\qq}[1]{\left[#1\right]_v}
\newcommand{\FF}{\mathbf{F}}

\title{}
\author{}
\date{}

\begin{document}
\section*{\texttt{tl\_jones\_ff.py}}
Code to compute Jones representation trace polynomials $\mathrm{tr}(\rho_\lambda(\beta))$, often more efficiently than symbolic matrix multiplication, using finite-field evaluations. 

\subsection*{Implementation}
We compute the trace polynomial $\mathrm{tr}(\rho_\lambda(\beta))\in\mathbb{Z}[v,v^{-1}]$ as follows:
\begin{enumerate}
\item Write $\beta=\sigma_{i_1}^{\varepsilon_1}\cdots\sigma_{i_L}^{\varepsilon_L}$ with $\varepsilon_j \in\{\pm1\}$ and choose a large prime $p$ and a primitive root of unity $\omega$ in $\mathbb{F}_p^\times$. 
\item Evaluate the Jones representation matrices of each braid generator at $v=\omega^0,\omega^1,\ldots,\omega^{2L}$,
\item Multiply these evaluated generator matrices to form evaluated braid matrices,
\item Take traces of these evaluated braid matrices to form evaluated traces in $\mathbb{F}_p$,
\item Use inverse DFT on these evaluated traces to form the trace polynomial's reduction mod $p$,
\item Lift the trace polynomial's reduction mod $p$ to form the true trace polynomial. 
\end{enumerate}

We explain the choice of $p$ and $\omega$ now. We know $\mathrm{tr}(\rho_\lambda(\beta))$ is of the form
\[
t(v):=\sum_{m=-L}^{L} c_m v^m.
\]
So it suffices to compute the shifted polynomial
\[
s(v):=v^L\, t(v) = \sum_{j=0}^{2L} a_j v^j,
\qquad a_j=c_{j-L}.
\]
These coefficients are determined by \(2L+1\) evaluations. Suppose the input word is nonempty, so that $L>0$. Define $C := d^L$, where $d:= \dim(V_\lambda)$. Choose the smallest prime \(p\) such that 
\[
p>2C
\qquad
\text{and}
\qquad
p\equiv 1\ \text{mod}\ {2L+1}.
\]
Furthermore, choose a primitive $(2L+1)^\text{th}$ root of unity $\omega$ in $\mathbb{F}_p$. 

The latter condition on $p$ ensures the root $\omega$ exists. The former condition on $p$ ensures we can determine each coefficient $c_m$, and therefore each shifted coefficient $a_j$, from its residue modulo $p$. This is because each coefficient lies in the interval $[-C,C]$, which has length $2C<p$.

Now suppose we have evaluated traces $t(\omega^k) \in\mathbb{F}_p$ for $k=0,1,\ldots,2L$. Then $s(\omega^k)=\omega^{kL}\, t(\omega^k)$ and we have
\[
\omega^{kL}\, t(\omega^k) = \sum_{i=0}^{2L} a_i\omega^{ki}.
\]
The inverse discrete Fourier transform gives
\[
a_j
=\frac{1}{2L+1}\sum_{k=0}^{2L}
\omega^{kL}t(\omega^k)\,\omega^{-kj}
\quad\text{in }\mathbb{F}_p,
\]
which recovers the shifted coefficients modulo \(p\). Each residue is then lifted to its symmetric integer representative in the interval \([-p/2,p/2]\). Since the true coefficients lie in \([-C,C]\) and \(p>2C\), this lift gives the exact integer coefficient. Finally, the code unshifts by using \(c_{j-L}=a_j\), so
\[
t(v)=v^{-L}s(v)=\sum_{j=0}^{2L} a_j v^{j-L}.
\]


\newpage
\subsection*{Overview}


\subsubsection*{\texttt{tl\_jones\_ff.py}}
\begin{enumerate}
  \item \texttt{FFMatrix}
  \item \texttt{FFEvaluationContext}
  \item \texttt{ff\_jones\_polynomial}
  \item \texttt{ff\_jones\_rep\_trace}
  \item \texttt{ff\_jones\_rep\_braid\_word}
  \item \texttt{\_ff\_create\_context}
  \item \texttt{\_inverse\_dft}
  \item \texttt{\_ff\_temperley\_lieb\_matrix}
  \item \texttt{\_identity\_matrix}
  \item \texttt{\_matrix\_multiply\_mod}
\end{enumerate}

The main public entry point is \texttt{ff\_jones\_polynomial}, which computes the same unnormalised Jones polynomial as \texttt{tl\_jones.jones\_polynomial}, but computes each representation trace using finite-field interpolation.
The function \texttt{ff\_jones\_rep\_trace} computes one such trace polynomial.
The function \texttt{ff\_jones\_rep\_braid\_word} computes the evaluated braid matrices which feed into the trace computation.
The remaining functions are internal helpers for building the finite-field context, applying inverse DFT, representing sparse Temperley--Lieb action data, and multiplying dense finite-field matrices.


\newpage
\subsection*{Function Explanations}

\paragraph{\texttt{FFMatrix}}
This is a type alias for a dense finite-field matrix, represented as a list of rows, where each row is a list of Python integers.
The entries are always interpreted modulo the prime stored in the relevant \texttt{FFEvaluationContext}.

\paragraph{\texttt{FFEvaluationContext}}
Stores the shared finite-field data needed to evaluate and reconstruct one trace.
It contains the strand number \(n\), the index \(r\) for the module \(V_{(n-r,r)}\), the validated braid word, the prime \(p\), the interpolation order, the chosen root of unity, the degree bound, the coefficient bound, and the dimension of the link-state basis.

\paragraph{\texttt{ff\_jones\_polynomial}}
Computes the unnormalised Jones polynomial of the closure of a braid word.
It validates the braid word, computes its writhe, then forms the same weighted trace sum as \texttt{tl\_jones.jones\_polynomial}:
\[
(-v^2)^{-w(\beta)}
\sum_{r=0}^{\lfloor n/2\rfloor}
\qq{n-2r+1}\,
\Tr(\rho_{(n-r,r)}(\beta)).
\]
The difference is that each trace \(\Tr(\rho_{(n-r,r)}(\beta))\) is computed by \texttt{ff\_jones\_rep\_trace}.

\paragraph{\texttt{ff\_jones\_rep\_trace}}
Computes the Laurent polynomial \(\Tr(\rho_{(n-r,r)}(\beta))\) by finite-field interpolation.
It first builds an \texttt{FFEvaluationContext}, calls \texttt{ff\_jones\_rep\_braid\_word} to obtain the braid matrix evaluated at each interpolation point, takes the trace of each evaluated matrix modulo \(p\), applies \texttt{\_inverse\_dft}, and finally lifts the residues to symmetric integer representatives to recover the exact Laurent polynomial.

\paragraph{\texttt{ff\_jones\_rep\_braid\_word}}
Evaluates the Jones representation of a whole braid word at every interpolation root \(1,\omega,\ldots,\omega^{M-1}\), where \(M\) is the interpolation order.
For each root it builds each generator matrix from cached Temperley--Lieb action data, multiplies the generator matrices in braid-word order, and returns the resulting list of dense finite-field braid matrices.
This is the finite-field analogue of \texttt{tl\_jones.jones\_rep\_braid\_word}, except that it returns one evaluated matrix for each interpolation point rather than one symbolic matrix.

\paragraph{\texttt{\_ff\_create\_context}}
Validates \(n\), \(r\), and the braid word, computes \(d=\dim(V_{(n-r,r)})\), and chooses the prime and root of unity used for reconstruction.
For the empty word it uses degree bound \(0\), interpolation order \(1\), and root of unity \(1\).
For a nonempty word of length \(L\), it uses degree bound \(L\), interpolation order \(2L+1\), coefficient bound \(d^L\), and chooses a prime \(p>2d^L\) with \(p\equiv 1 \pmod{2L+1}\), so that a primitive \((2L+1)\)-st root of unity exists in \(\FF_p\).

\paragraph{\texttt{\_inverse\_dft}}
Recovers the shifted Laurent coefficients modulo \(p\) from the evaluated trace values.
Given values \(t(\omega^k)\), it first multiplies by \(\omega^{kL}\) to shift the Laurent polynomial into an ordinary polynomial of degree at most \(2L\), then applies the inverse discrete Fourier transform over \(\FF_p\).
The returned list contains residues corresponding to the powers \(v^{-L},v^{-L+1},\ldots,v^L\).

\paragraph{\texttt{\_ff\_temperley\_lieb\_matrix}}
Computes cached sparse column-action data for the Temperley--Lieb generator \(E_i\) on the link-state basis of \(V_{(n-r,r)}\).
For each basis vector it records the coefficient of the image and the row index of the resulting basis vector.
This stores the same information as \texttt{tl\_jones.\_temperley\_lieb\_matrix}, but avoids constructing a symbolic SymPy matrix.

\paragraph{\texttt{\_identity\_matrix}}
Builds the dense identity matrix of a given size as a list of lists of integers.
It is used as the starting matrix when multiplying the evaluated generator matrices for a braid word.

\paragraph{\texttt{\_matrix\_multiply\_mod}}
Multiplies two dense integer matrices modulo \(p\).
The implementation skips zero entries in the left and right factors, which is useful because many evaluated Temperley--Lieb generator matrices are sparse or nearly sparse.










\end{document}
